\par\Needspace{6\baselineskip}\medskip
\hypertarget{native-217-article158}{}\pdfbookmark[2]{158.}{article-158}
\begin{center}158.\end{center}
Duo exempla praecedentia occasionem nondum suppeditaverunt, methodum art.~120 in usum vocandi: hypotheses enim successivae tam rapide convergebant, ut iam in secunda subsistere licuisset, tertiaque a veritate vix sensibiliter aberraret. Revera hocce commodo semper fruemur, quartaque hypothesi supersedere poterimus, quoties motus heliocentricus modicus est, tresque radii vectores non nimis inaequales sunt, praesertim si insuper temporum intervalla parum inter se discrepant. Quanto magis autem problematis conditiones hinc recedunt, tanto fortius valores primi suppositi quantitatum $P$, $Q$ a veris different, tantoque lentius valores sequentes ad veros convergent. In tali itaque casu tres quidem primae hypotheses ita absolvendae sunt, uti duo exempla praecedentia monstrant (ea sola differentia, quod in hypothesi tertia non elementa ipsa, sed, perinde ut in hypothesi prima et secunda,
\SourceLeaf{218}
quantitates $\eta$, $\eta''$, $P'$, $Q'$, $X$, $Y$ computare oportet): dein vero haud amplius valores postremi ipsarum $P'$, $Q'$ tamquam valores novi quantitatum $P$, $Q$ in hypothesi quarta accipientur, sed hi per methodum art.~120 e combinatione trium primarum hypothesium eruentur. Rarissime tunc opus erit, ad hypothesin quintam secundum praecepta art.~121 progredi. --- Iam hos quoque calculos exemplo illustrabimus, ex quo simul elucebit, quam late methodus nostra pateat.
